% !TEX root = ../main.tex
\section{Elementary Properties of PageRank and Rank Correlation}
\label{ch:appendix-theory}

This appendix records elementary properties used implicitly in Chapters~\ref{ch:literature}--\ref{ch:results}. The material is standard \cend{49,40,55,36} and is included for self-contained reading.

\dissertationSubheading[sec:pr-existence]{Existence and uniqueness of PageRank}

Let $P$ be a row-stochastic matrix on $n$ nodes, possibly after dangling-node repair as in Equation~\eqref{eq:transition}. For $d\in(0,1)$, the Google matrix
\[
G_{\mathrm{PR}}=d P^{\top}+\frac{1-d}{n}\mathbf{1}\mathbf{1}^{\top}
\]
is positive and column-stochastic (up to the transpose convention in Equation~\eqref{eq:google-matrix}). By the Perron--Frobenius theorem for positive matrices, $G_{\mathrm{PR}}$ has a unique eigenvalue of modulus one with a strictly positive eigenvector $\boldsymbol{\pi}$ normalized so that $\sum_i\pi_i=1$. Hence PageRank scores exist and are unique for any directed graph once damping and dangling-node rules are fixed \cend{40}.

This uniqueness underpins comparisons between EndorseRank and AWP: given fixed parameters, each method induces a single score vector on a given edge set. Differences between $\mathbf{s}_{\mathrm{ER}}$ and $\mathbf{s}_{\mathrm{AWP}}$ are therefore attributable to edge semantics and weights, not to solver multiplicity.

\dissertationSubheading[sec:pr-continuity]{Continuity in edge weights}

PageRank is a continuous function of the entries of $P$ for fixed $d\in(0,1)$. Small perturbations of allowance or transfer weights therefore induce small perturbations of scores in the $\|\cdot\|_1$ norm, with Lipschitz constants depending on $d$ and $n$ \cend{40}. Continuity justifies robustness checks that vary damping or restrict to top tokens: if alignment were knife-edge fragile, small hyperparameter changes would scramble family-level $\tau$. Empirically, EndorseRank allowance $\tau$ remains $\approx 0.355$ across $d\in\{0.75,0.85,0.95\}$ (Table~\ref{tab:robustness-damping}), consistent with a stable ranking geometry.

\dissertationSubheading[sec:pr-sparsity]{Sparsity and per-iteration cost}

A sparse matrix--vector product $P^{\top}\mathbf{x}$ costs $O(m)$ arithmetic operations for $m$ nonzero edges. Power iteration therefore costs $O(mI)$ for $I$ iterations. On the matched cohort, $m_{\mathrm{ER}}=14{,}727$ and $m_{\mathrm{AWP}}=137{,}087$, so even at identical $I$ the transfer method would be roughly $9.3\times$ more expensive per iteration. Observed iteration counts ($I_{\mathrm{ER}}\approx 37$, $I_{\mathrm{AWP}}\approx 100$) widen the gap, matching the measured wall-time ratio $\approx 8.9\times$ (Table~\ref{tab:benchmark-runtime}).

\dissertationSubheading[sec:spearman-properties]{Properties of Spearman's $\rho$}

Spearman's $\rho$ is Pearson's correlation applied to ranks. It is invariant to strictly increasing transformations of either variable and lies in $[-1,1]$. Perfect monotone agreement yields $\rho=1$; perfect reverse agreement yields $\rho=-1$. Ties require mid-rank adjustments; the evaluation pipeline uses standard tie-aware implementations. Because reputation scores and proxies are continuous and heavy-tailed, rank-based association is more appropriate than Pearson correlation on raw values \cend{16}.

\dissertationSubheading[sec:kendall-properties]{Properties of Kendall's $\tau$}

Kendall's $\tau$ equals the normalized difference between concordant and discordant pairwise orderings. It admits a probabilistic reading: for two random distinct wallets, $\tau$ is the probability their relative order agrees under both rankings, minus the probability it disagrees \cend{36}. This reading is useful when communicating results to practitioners: $\tau=0.355$ means a clear but imperfect tendency for allowance-proxy orderings to agree with EndorseRank orderings.

For large $n$, $\tau$ and $\rho$ are often similar in sign and magnitude but not identical. This dissertation reports both, following Do et al.\cend{16}, and uses family-mean $\tau$ as the primary summary.

\dissertationSubheading[sec:construct-validity-note]{Construct validity note}

Cronbach and Meehl\cend{14} distinguish criterion-related validity from construct validity. In the absence of a gold-standard default label, this dissertation pursues a construct-oriented strategy: each method is tested against proxies that operationalize its intended meaning (allowance receipt for EndorseRank; inbound transfer activity for AWP) and against secondary families that probe neighboring constructs (risk, stability, trading success). High within-construct $\tau$ supports interpretation; low cross-construct $\tau$ is not automatically failure---it may indicate discriminant validity. GF-PR's extreme trading-success $\tau$ is a cautionary case of \emph{contaminated} criteria, where the criterion shares inputs with the measure (SRQ1).

\dissertationSubheading[sec:numerical-summary-sheet]{Numerical summary sheet}

For quick reference, the primary matched-cohort numerics are:

\begin{center}
\begin{tabular}{ll}
\toprule
Quantity & Value \\
\midrule
Matched wallets $n$ & $5{,}521$ \\
GMX closes (decoded) & $365{,}488$ \\
EndorseRank edges & $14{,}727$ \\
AWP edges & $137{,}087$ \\
EndorseRank runtime & $2.105$\,s \\
AWP runtime & $18.711$\,s \\
Runtime ratio (AWP/ER) & $\approx 8.9\times$ \\
Inter-method $\tau$ & $0.316$ \\
ER allowance mean $\tau$ & $0.355$ \\
AWP transfer mean $\tau$ & $0.534$ \\
ER trading-success mean $\tau$ & $0.096$ \\
AWP trading-success mean $\tau$ & $0.179$ \\
GF-PR trading-success mean $\tau$ & $0.583$ \\
GF-PR inverse-risk mean $\tau$ & $0.376$ \\
Expanded pool $N$ & $31{,}612$ \\
Speedup at $n=10{,}000$ & $\approx 10.4\times$ \\
Speedup at $N=31{,}612$ & $\approx 9.4\times$ \\
Observation window & 2025-12-01 -- 2026-05-31 \\
\bottomrule
\end{tabular}
\end{center}

These values are authoritative only insofar as they match \texttt{eval\_summary.json} and the exported LaTeX tables; if a regenerated export differs, the JSON artifact prevails.
